%!TEX encoding = UTF8
%!TEX root = 0-notes.tex

\unchapter{Nombres rationnels}
\fancyhead[C]{\textbf{Chapitre \thechapter~— Nombres rationnels}}
\addcontentsline{toc}{section}{Méthodologie}


\dfn{nombres rationnels}{
	Soit $a, b \in \Z$ deux entiers relatifs, avec $b$ non nul.
	
	Alors le \emphindex{rationnel} $\frac{a}{b}$ est l'\underline{unique} nombre vérifiant
		\[ \frac{a}b \times b = a. \]
	
	L'ensemble de tous les rationnels est noté
		\[ \Q = \Bigset{ \frac{a}{b} \tq a \in \Z, b \in \Z, b \neq 0}. \]
}{dfn:nombres-rationnels}

\ex{}{
	Le nombre rationnel $\frac43$ est l'unique nombre vérifiant
		\[ \frac43 \times 3 = 4. \]
}{ex:nombres-rationnels2}

\thm{opérations}{
	Pour $a, b, c, d \in \Z$ des entiers relatifs, non nuls lorsque dénominateurs,
		\begin{align*}
		\frac{a}b + \frac{c}d = \frac{ad + cb}{bd}, && \et &&
		\frac{a}b \times \frac{c}d = \frac{a\times c}{b\times d}.
		\end{align*}
}{}

\dfn{}{
	L'\emphindex{inverse} d'un nombre $a \in \R$ non nul désigne la quantité qui, lorsque multipliée par $a$, donne $1$.
	Elle est notée $\frac{1}{a}$ ou $a^{-1}$ et vérifie 
		\[ a \times \frac1a = \frac{a}{a} = 1. \]
	
	Phrase à retenir :
	\begin{center}
		« Diviser par $a$ c'est multipler par son inverse $\frac1a$. »
	\end{center}
}{}

\nt{
	L'inverse de $2$ est donc $0,5$ car $2\times0,5 = 1$. 
	D'où $\frac12 = 0,5$.
	
	La notation $\frac12$ est préférée à $0,5$ car généralement plus précise : le rationnel $\frac13$ n'est pas décimal et ne peut pas s'exprimer de façon exacte en écriture décimale (elle est infinie).
}
	
\ex{}{
	L'inverse de $\frac47$ est $\frac74$ car $\frac47 \times \frac74 = 1$.
	Autrement dit, $\frac{\ \ 1 \ \ }{\frac47} = \left( \frac47 \right)^{-1} = \frac74$.
}{}

\dfn{fraction irréductible}{
	La relation suivante permet de réduire les fractions.
		\[ \frac{a\times d}{b\times d} = \frac{a}{b} \]
	Lorsqu'un diviseur commun au numérateur et au dénominateur est trouvé, il peut être simplifié.
	Sinon la fraction est \emphindex{irréductible}.
}{}

\thm{}{
	Les nombres $\sqrt2, \oldsqrt[3]{7}, \pi$ ne sont pas rationnels.
	
	Ils appartiennent à $\R$, l'ensemble des \emphindex{nombres réels}.
}{}

\newpage
\addcontentsline{toc}{section}{Exercices}
\fancyhead[C]{\textbf{Chapitre \thechapter~— Exercices}}


\exe{}{
	Écrire les sommes et différences de rationnels sous forme d'une fraction irréductible.
	\begin{multicols}{3}
	\begin{enumerate}[itemsep=10pt]
		\item $1 - \dfrac12$
		\item $\dfrac78 + \dfrac{9}{40}$
		\item $\dfrac{3}{4} + \dfrac{2}{5}$
		\item $\dfrac{7}{6} - \dfrac{5}{8}$
		\item $\dfrac{9}{10} + \dfrac{2}{3} - \dfrac{1}{4}$
		\item $\dfrac{-5}{7} + \dfrac{8}{9}$
%		\item $\dfrac{4}{5} - \dfrac{3}{10} + \dfrac{-5}{6}$
%		\item $\dfrac{-1}{2} + \dfrac{7}{8}$
%		\item $\dfrac{11}{12} - \dfrac{5}{9} + \dfrac{7}{10}$
%		\item $\dfrac{2}{3} - \dfrac{5}{4}$
%		\item $\dfrac{-1}{3} + \dfrac{7}{9} - \dfrac{5}{6}$
%		\item $\dfrac{3}{8} - \dfrac{1}{2}$
	\end{enumerate}
	\end{multicols}
}{exe:rationnals}
{
	\begin{multicols}{2}
	\begin{enumerate}
		\item $1 - \dfrac{1}{2} = \dfrac{1}{2}$
		\item $\dfrac{7}{8} + \dfrac{9}{40} = \dfrac{11}{10}$
		\item $\dfrac{3}{4} + \dfrac{2}{5} = \dfrac{23}{20}$
		\item $\dfrac{7}{6} - \dfrac{5}{8} = \dfrac{13}{24}$
		\item $\dfrac{9}{10} + \dfrac{2}{3} - \dfrac{1}{4} = \dfrac{79}{60}$
		\item $\dfrac{-5}{7} + \dfrac{8}{9} = \dfrac{11}{63}$
%		\item $\dfrac{4}{5} - \dfrac{3}{10} + \dfrac{-5}{6} = \dfrac{-1}{3}$
%		\item $\dfrac{-1}{2} + \dfrac{7}{8} = \dfrac{3}{8}$
%		\item $\dfrac{11}{12} - \dfrac{5}{9} + \dfrac{7}{10} = \dfrac{191}{180}$
%		\item $\dfrac{2}{3} - \dfrac{5}{4} = \dfrac{-7}{12}$
%		\item $\dfrac{-1}{3} + \dfrac{7}{9} - \dfrac{5}{6} = \dfrac{-7}{18}$
%		\item $\dfrac{3}{8} - \dfrac{1}{2} = \dfrac{-1}{8}$
	\end{enumerate}
	\end{multicols}
}


\exe{}{
	Écrire les produits et divisions de rationnels sous forme de fraction irréductible.
	
	\begin{multicols}{3}
	\begin{enumerate}[itemsep=10pt]
		\item $\dfrac{\ \ -3 \ \ }{\dfrac54}$
		\item $\dfrac{63}{20} \times \dfrac{15}{14}$
		\item $\dfrac{3}{4} \times \dfrac{2}{5}$
		\item $\dfrac{\ \ \dfrac{7}{6}\ \ }{\dfrac{5}{8}}$
		\item $\dfrac{9}{10} \times \dfrac{2}{3} \times \dfrac{4}{1}$
		\item $\dfrac{-5}{7} \times \dfrac{8}{9}$
%		\item $\left(\dfrac{4}{5}\right)^{-1} \times \dfrac{3}{10}$
%		\item $\dfrac{-1}{2} \times \left(\dfrac{7}{8}\right)^{-1}$
%		\item $\dfrac{11}{12} \times \dfrac{5}{9} \times \left(\dfrac{7}{10}\right)^{-1}$
%		\item $\dfrac{2}{3} \times \dfrac{5}{4}$
%		\item $\dfrac{-1}{3} \times \left(\dfrac{7}{9}\right)^{-1}$
%		\item $\dfrac{3}{8} \times \dfrac{1}{2}$
	\end{enumerate}
	\end{multicols}
}{exe:rationnals2}
{
	\begin{multicols}{2}
	\begin{enumerate}[itemsep=10pt]
		\item $\dfrac{-3}{\dfrac{5}{4}} = \dfrac{-12}{5}$
		\item $\dfrac{63}{20} \times \dfrac{15}{14} = \dfrac{27}{8}$
		\item $\dfrac{3}{4} \times \dfrac{2}{5} = \dfrac{3}{10}$
		\item $\dfrac{\ \ \dfrac{7}{6}\ \ }{\dfrac{5}{8}} = \dfrac{28}{15}$
		\item $\dfrac{9}{10} \times \dfrac{2}{3} \times \dfrac{4}{1} = \dfrac{12}{5}$
		\item $\dfrac{-5}{7} \times \dfrac{8}{9} = \dfrac{-40}{63}$
%		\item $\left(\dfrac{4}{5}\right)^{-1} \times \dfrac{3}{10} = \dfrac{3}{8}$
%		\item $\dfrac{-1}{2} \times \left(\dfrac{7}{8}\right)^{-1} = \dfrac{-4}{7}$
%		\item $\dfrac{11}{12} \times \dfrac{5}{9} \times \left(\dfrac{7}{10}\right)^{-1} = \dfrac{275}{378}$
%		\item $\dfrac{2}{3} \times \dfrac{5}{4} = \dfrac{5}{6}$
%		\item $\dfrac{-1}{3} \times \left(\dfrac{7}{9}\right)^{-1} = \dfrac{-3}{7}$
%		\item $\dfrac{3}{8} \times \dfrac{1}{2} = \dfrac{3}{16}$
	\end{enumerate}
	\end{multicols}
}


\exe{}{
	Exprimer chaque valeur suivante sous forme de fraction irréductible.
	\begin{multicols}{4}
	\begin{enumerate}[label=\alph*), itemsep=10pt]
		\item $1-\dfrac{4}{44}$
		\item $\dfrac{15}{27}$
		\item $\dfrac{8}{3} + 6$
		\item $\dfrac{65}{5^2}$
		\item $\dfrac{6}{32}\times\dfrac{2}{3}$
		\item $\dfrac{38}{20}-3$
		\item $\left(\dfrac{10}{44}\right)^{-1}$
		\item $\dfrac{6}{18}\times2 + 1$
		\item $\left(\dfrac{66}{99}\right)^2$
		\item $\dfrac{23}{18} + \dfrac12$
		\item $\dfrac{34}{20}$
		\item $2 - \dfrac{9}{5}$
	\end{enumerate}
	\end{multicols}
}{exe:frac-irr}{
	\begin{multicols}{4}
	\begin{enumerate}[label=\alph*), itemsep=10pt]
		\item $\dfrac{10}{11}$
		\item $\dfrac{5}{9}$
		\item $\dfrac{26}{3}$
		\item $\dfrac{13}{5}$
		\item $\dfrac{1}{18}$
		\item $\dfrac{-11}{10}$
		\item $\dfrac{22}{5}$
		\item $\dfrac{5}{3}$
		\item $\dfrac{4}{9}$
		\item $\dfrac{16}{9}$
		\item $\dfrac{17}{10}$
		\item $\dfrac{1}{5}$
	\end{enumerate}
	\end{multicols}
}


\exe{}{
	Sans calculatrice, donner le développement décimal des fractions suivantes.
	\begin{multicols}{3}
	\begin{enumerate}[label=\roman*), leftmargin=60pt]
		\item $\dfrac1{10}$
		\item $\dfrac1{5}$
		\item $\dfrac3{10^{5}}$
		\item $\dfrac7{20}$
		\item $\dfrac{395}{50}$
		\item $\dfrac{11}{200}$
	\end{enumerate}
	\end{multicols}
}{exe:dev-decimaux}{
	\begin{multicols}{3}
	\begin{enumerate}[label=\roman*)]
		\item $\dfrac1{10} = 0,1.$
		\item $\dfrac1{5} = \dfrac2{10} = 0,2.$
		\item $\dfrac3{10^{5}} = 0,000~03.$
		\item $\dfrac7{20} = \dfrac{3,5}{10} = 0,35.$
		\item $\dfrac{395}{50} = \dfrac{790}{100} = 7,9.$
		\item $\dfrac{11}{200} = \dfrac{5,5}{100} = 0,055.$
	\end{enumerate}
	\end{multicols}
}


\exe{}{
	Sans calculatrice, donner le développement décimal des fractions suivantes.
	\begin{multicols}{3}
	\begin{enumerate}[label=\roman*), leftmargin=60pt]
		\item $\dfrac1{100}$
		\item $\dfrac3{5}$
		\item $\dfrac2{10^{1}}$
		\item $\dfrac7{2}$
		\item $\dfrac{7}{20}$
		\item $\dfrac{31}{200}$
	\end{enumerate}
	\end{multicols}
}{exe:dev-decimauxE}{
	\begin{multicols}{3}
	\begin{enumerate}[label=\roman*)]
		\item $\dfrac1{100} = 0,01$
		\item $\dfrac3{5} = 0,6$
		\item $\dfrac2{10^{1}} = 0,2$
		\item $\dfrac7{2} = 3,5$
		\item $\dfrac{7}{20} = 0,35$
		\item $\dfrac{31}{200} = 0,155$
	\end{enumerate}
	\end{multicols}
}


\exe{, difficulty=1}{
	Déterminer le développement décimal de $\frac1{11}$ sans calculatrice.
}{exe:111-dec}{
	Posons 
		\[ \dfrac{1}{11} = 0,n_1 n_2 n_3 n_4 \dots. \]
	Par multplication par 10 successives et étude des parties décimales, on obtient $n_1 = 0$, $n_2=9$, et $\frac1{11} = 0,n_3n_4n_5\dots$.
	D'où
		\[ \dfrac1{11} = 0,09~09~09~09~\cdots. \]
}


\exe{, difficulty=1}{
	Déterminer le développement décimal de $\frac{12}{22}$ sans calculatrice.
}{exe:113-dec}{
	Posons 
		\[ \dfrac{12}{22} = 0,n_1 n_2 n_3 n_4 \dots. \]
	Par multplication par 10 successives et étude des parties décimales, on obtient $n_1 = 5$, $n_2=4$, et $\frac{12}{22} = 0,n_3 n_4 n_5\dots$.
	D'où
		\[ \dfrac{12}{22} = 0,54~54~54~54~\cdots. \]
}


\exe{}{
	Montrer que les nombres suivants sont rationnels en les exprimant sous forme de fraction d'entiers.
	\[
	\begin{aligned}
		A &= 0,666{\dots} \\
		B &= 1,666{\dots} \\
	\end{aligned}
	\hspace{5cm}
	\begin{aligned}
		C &= 0,121212{\dots} \\
		D &= 0,34777{\dots} \\
	\end{aligned}
	\]
	\[
	E = 0,123456789123456789123{\dots} \text{ (nombre d'Amandine) }
	\]
}{exe:dev-to-fraction}{
	\begin{enumerate}[label=\Alph*.]
		\item 
		La relation $10A = 6 + A$ donne $A = \frac23$.
		\item
		La relation $10B = 15 + B$ donne $B = \frac{15}{9} = \frac53$.
		On aurait aussi pû utiliser que $B = 1+A$.
		\item
		La relation $100C = 12 + C$ donne $C = \frac{12}{99} = \frac{4}{33}$.
		\item
		En posant $\tilde{E} = 100E = 34,777{\dots}$, on obtient $10\tilde{E} = 313 + \tilde{E}$, d'où $\tilde{E} = \frac{313}{9}$, et donc $E = \frac{1}{100}E' = \frac{313}{900}$.
		\item 
		La relation $10^9 E = 123~456~789 + E$ donne $E = \frac{123~456~789}{999~999~999}$.
	\end{enumerate}
}


\exe{}{
	Montrer que les nombres suivants sont rationnels en les exprimant sous forme de fraction d'entiers.
	\begin{align*}
		A = 0,123123123{\dots} &&
		B = 0,9999{\dots}
	\end{align*}
}{exe:dev-to-fractionE}{
	\begin{enumerate}[label=\Alph*.]
		\item
		La relation $1000A = 123 + A$ donne $A = \frac{123}{999} = \frac{41}{333}$.
		\item 
		La relation $10B = 9 + B$ donne $9B = 9$ et $B=1$.
	\end{enumerate}
}


\exe{, difficulty=2}{
	Donner quelques termes en plus de la somme infinie suivante. 
		\[ \dfrac1{10} + \dfrac1{100} + \dfrac1{1 000} + \dfrac1{10 000} + \dots \]
	Écrire le nombre en écriture décimale, puis l'exprimer sous forme de fraction de deux entiers comme à l'exercice \ref{exe:dev-to-fraction}.
}{exe:20}{
	La somme infinie vaut $0,1111{\dots} = \frac19$.
}

\exe{, difficulty=2}{
	Donner quelques termes en plus de la somme infinie suivante.
		\[ \dfrac12 + \dfrac14 + \dfrac18 + \dfrac1{16} + \dots \]
	À quoi est-elle égale ? Un résultat entier est attendu.
}{exe:21}{
	Posons $S$ cette somme. 
	En multipliant par $2$, on trouve $2S = 1 + S$, et donc $S=1$.
}

\exe{, difficulty=2}{
	On considère la somme finie suivante
		\[ S = 1 + 2 + 4 + 8 + 16 + \dots + 2^{n-1} +  2^n, \]
	où $n\in\N$ est un entier naturel.
	
	Montrer que $2S = S - 1 + 2^{n+1}$, et en déduire la valeur de $S$ sous forme close (sans points de suspension).
}{exe:22}{
	En multipliant par 2, le premier terme (1) disparait, et un nouveau terme ($2^{n+1}$) apparaît.
	On a donc bien $2S = S -1 + 2^{n+1}$, et donc $S = 2^{n-1} - 1$.
}


\exe{, difficulty=1}{
	Soit $q\in\Q$ un rationnel. Montrer qu'un triangle dont les côtés sont de longueur $(4 ; 5 ; q)$ ne peut être rectangle que si $q=3$.
	
	En autorisant $x\in\R$ irrationel, donner un triangle rectangle de côtés $(4 ; 5; x)$ avec $x\neq3$.
}{exe:non-rec}{
	Il y a deux cas de figure à étudier : $q$ est la longueur de l'hypothénuse, ou elle ne l'est pas.
	Dans le premier cas, le théorème de Pythagore énonce que
		\[ 4^2 + 5^2 = q^2 \iff 41 = q^2. \]
	Il suit que $q=\sqrt{41}$ qui n'est pas rationnel car 41 n'est pas un carré parfait.
	Ceci répond à la deuxième question : $(4 ; 5; \sqrt{41})$ est rectangle.
	
	Dans le second cas,
		\[ q^2 + 4^2 = 5^2 \iff q^2 = 9. \]
	Par conséquent, $q = \pm3$ et seule la solution $q=3$ est positive et donc une longueur valide.
}

\exe{, difficulty=2}{
	Montrer que $\sqrt2 = 1 + \frac1{1+\sqrt2}$.
	En déduire que $\sqrt2 = 1 + \frac1{2 + \frac1{1+\sqrt2}}$ et que
		\[\sqrt2 = 1 + \cfrac1{2 + \cfrac1{2 + \cfrac1{2 + \cfrac1{1+\sqrt2}}}}. \]
	Exprimer le rationnel $1 + \frac1{2 + \frac1{2 + \frac12}}$ sous la form $\frac{p}{q}$ avec $p$ et $q$ premiers entre eux et calculer $\left| \sqrt2 - \frac{p}q \right|$ à l'aide de la calculatrice.
	Faire idem avec $1 + \frac1{2 + \frac1{2 + \frac1{2 + \frac12}}}$.
}{exe:frac-continue-sqrt2}{
	En manipulant l'équation on se retrouve à vouloir montrer que $\sqrt2 - 1 = \frac1{\sqrt2+1}$, ce qui est équivalent à $(\sqrt2 - 1)(\sqrt2 + 1) = 1$, vrai par identité remarquable (ou par distributivité simple).
	
	Comme $\sqrt2 = 1 + \frac1{1+\sqrt2}$, on peut remplacer le $\sqrt2$ dans l'expression de droite par toute l'expression de droite.
	On obtient donc $\sqrt2 = 1 + \cfrac1{1 + 1 + \cfrac1{1+\sqrt2}} = 1 + \cfrac1{2 + \cfrac1{1+\sqrt2}}$ comme voulu.
	On obtient l'expression d'après en répétant le même raisonnement.
	
	Par simplifications simples, on trouve $\frac{p}{q} = \frac{17}{12}$.
	La calculatrice nous donne alors $\left| \sqrt2 - \frac{17}{12} \right| \approx 0,002$.
	
	Finalement, remarquons que 
	\begin{align*}
		1 + \cfrac1{2 + \cfrac1{2 + \cfrac1{2 + \cfrac12}}} &= 1 + \cfrac1{1+ 1 + \cfrac1{2 + \cfrac1{2 + \cfrac12}}}, \\
			&= 1 + \cfrac1{1+ \dfrac{17}{12}} = \dfrac{41}{29}.
	\end{align*}
	Le rationnel $\frac{41}{29}$ vérifie $\left| \sqrt2 - \frac{41}{29} \right| \approx 0,0004$, soit 5 fois plus proche de $\sqrt2$ que $\dfrac{17}{12}$ l'est.
}

\exe{, difficulty=2}{
	Montrer qu'un triangle isocèle rectangle ne peut pas avoir des côtés de longueurs toutes entières.
}{exe:non-iso-rec}{
	Supposons qu'il existe un triangle isocèle rectangle de côtés $(a ; a; b)$.
	
	Or, comme le triangle est rectangle isocèle,
		\[ b^2 = a^2 + a^2 = 2a^2 \iff b = \pm\sqrt2 a. \]
	Par conséquent, $\sqrt2 = \pm\frac{b}{a} \in \Q$, ce qui contredit le théroème d'irrationalité de la racine carrée de 2.
}

\exe{, difficulty=2}{
	En musique, la gamme classique contient douze notes avec un rapport de fréquence de 2 sur l'octave.
	La gamme est dite à tempérament égal si le rapport de fréquence entre deux notes consécutives est toujours le même.
	Notons ce rapport $r\in\R$.
	Justifier que $r^{12} = 2$ puis démontrer que $r$ est un nombre irrationnel : $r \not\in\Q$.
}{exe:temp-egal}{
	Le rapport de fréquence entre deux notes consécutives étant toujours $r$, on multiplie $r$ par lui-même 12 fois pour obtenir le rapport de fréquence entre deux notes de distance $12$, c'est-à-dire deux notes séparée par un octave.
	Ainsi $r^{12} = 2$ car l'octave correspond à un rapport de fréquence de 2.
	
	Supposons désormais que $r$ soit rationnel. Alors $r^{6}$ est aussi rationnel.
	Cependant, $(r^6)^2 = r^{12} = 2$, donc $r^{6} = \sqrt2$, qui n'est pas rationnel. \Large\Lightning
}

\exe{, difficulty=3}{
	Posons $\Qsqrt{5} = \bigset{ r + s \sqrt5 \tq r, s \in \Q} \subseteq \R$ et considérons $x = r + s\sqrt5$ et $y=p+q\sqrt5$ avec $r, s, p, q \in\Q$.
	\begin{enumerate}
		\item
		Montrer que $\Q \subseteq \Qsqrt5$ est une inclusion stricte.
		\item
		Montrer que $x=y \iff r = p \et s = q$.
	\end{enumerate}
	Considérons désormais la fonction « conjugué » $f$ définie par $f\bigl(r + s\sqrt5\bigr) = r - s\sqrt5$.
	\begin{enumerate}[resume]
		\item
		Montrer que $f(x+y) = f(x)+f(y)$ et que $f(xy) = f(x)f(y)$.
		\item
		Montrer que si $f(x) = x$, alors $x\in\Q$.
		\item
		Conclure à l'aide des deux dernières questions que
			\[ \dfrac1{2^{256}\sqrt5}\left[ \left({1+\sqrt5}\right)^{256} - \left({1-\sqrt5}\right)^{256}\right] \]
		est un nombre rationnel.
		(C'est en fait un nombre entier et le terme $256$ de la suite de Fibonacci)
		\item
		Montrer que $\Qsqrt{5} \neq \R$ en montrant que $\sqrt2 \not\in \Qsqrt{5}$.
	\end{enumerate}
} {exe:Qsqrt2}{
	%À rendre pour possibilité de points bonus à la prochaine évaluation.
	%Ni correction ni points ne seront attribués à un travail suspect ou ne démontrant pas une bonne compréhension de l'exercice.
	
	Indication pour la question 2 : il est possible de montrer en premier lieu que $x = 0 \iff r=s=0$.
	On dit alors que $1$ et $\sqrt5$ sont \emph{linéairement indépendants} sur $\Q$.
	
	Indication pour la question 5 : montrer que $f(x^n) = f(x)^n$ et que $f\left(\frac1x\right) = \frac1{f(x)}$ à l'aide de la question 3.
	
	\underline{Discussion}
	
	L'exercice fait l'étude du corps d'extension des nombres rationnels auxquels on a ajouté la racine carrée de 5.
	Il suit de la question 6 que l'ensemble 
		\[ \Q\bigl(\sqrt2, \sqrt5\bigr) = \bigset{ r + s \sqrt2 + t\sqrt5 \tq r, s, t \in \Q} \]
	est un sur-ensemble strict de $\Qsqrt5$ qui n'est toujours pas égal à $\R$.
	
	La {théorie de Galois}\footnotemark généralise ces exemples et permet notamment de répondre par la négative aux trois grands problèmes de l'Antiquité : la duplication du cube, la trisection de l'angle, et la quadrature du cercle.
	\footnotetext{Évariste Galois (1811-1832), mathématicien français mort à l'âge de 20 ans. }
	La démonstration procède ainsi: chaque construction à la règle et au compas à partir d'un repère orthonormé ajoute la racine carrée d'un nombre déjà constructible jusqu'ici.
	La racine carrée de 2 est par exemple constructible grâce au théorème de Pythagore comme distance entre $(1 ; 0)$ et $(0;1)$.
	En montrant que certains nombres ne peuvent pas appartenir à de telles extensions, on démontre immédiatement qu'ils sont non constructibles.
}


%%%%%%%%%%%%%%%%%%%%
%%%%% Samoyeau below %%%%%
%%%%%%%%%%%%%%%%%%%%



%\exe{difficulty=0}{
%	Simplifier les fractions suivantes quand c'est possible:
%	\\
%	$\frac{32}{28}$ ; $\frac{64}{80}$ ; $\frac{15}{35}$ ; $\frac{49}{35}$ ; $\frac{14}{21}$ ; $\frac{8}{16}$ ; $\frac{120}{140}$ ; $\frac{12}{36}$ ; $\frac{3700}{1200}$ ; $\frac{48}{56}$ ; $\frac{81}{99}$ ; $\frac{77}{66}$; $\frac{8}{12}$; $\frac{21}{27}$; $\frac{13}{45}$; $\frac{30}{75}$; $\frac{3}{2222}$; $\frac{12}{36}$.
%}{exe:Samoyeau1}{
%	sol1
%}

\exe{difficulty=0}{
	Classer les fractions suivantes par ordre croissant : 
	\begin{align*}
		A=\frac{1}{2}, && B=\frac{1}{6}, && C=\frac{5}{12}, &&D=\frac{3}{4}, &&E=\frac{2}{3}.
	\end{align*}
}{exe:Samoyeau2}{
	sol1
}

%\exe{difficulty=0}{
%	Calculer et simplifier:
%	\\
%	$A=\frac{2}{7}+\frac{6}{7}$ ;  $B=\frac{1}{3}-\frac{5}{3}$ ;  $C=\frac{5}{11}+\frac{-3}{11}$ ; $D=\frac{3}{8}+\frac{3}{4}$ ; $E=\frac{4}{9}+\frac{1}{27}$ ; $F=\frac{4}{30}+\frac{1}{10}$ ; $G=\frac{4}{5}+1$ ;$H=\frac{8}{3}-3$ ; $I=\frac{11}{13}+2$; $J=\frac{13}{15}-\frac{11}{20}$; $K=3+\frac{3}{4}$; $L=-1-\frac{5}{7}$; $M=\frac{1}{7}+\frac{3}{4}$; $N=\frac{5}{6}+\frac{-3}{8}$.
%
%}{exe:Samoyeau3}{
%	sol1
%}

%\exe{difficulty=0}{
%	Calculer et simplifier. 
%	\\
%	$A=\frac{7}{3} \times (-2)$ ; $B=\frac{9}{3} \times\frac{3}{8}$; $C=\frac{1}{5} \times \frac{-1}{2}$;$D=\frac{11}{15}\div \frac{-3}{2}$ ; $E=\frac{10}{3}\div \frac{1}{11}$; $F=\frac{\frac{2}{3}}{\frac{4}{9}} $, $G=\frac{\frac{5}{9}}{\frac{8}{3}} $, $H=\frac{\frac{3}{4}}{\frac{3}{4}} $, $I=\frac{\frac{3}{2}}{\frac{2}{3}} $, $J=\frac{11}{15}\times \frac{-3}{2}$.
%
%}{exe:Samoyeau4}{
%	sol1
%}

\exe{difficulty=0}{
	 Calculer et simplifier:
	\begin{align*}
	a&=\frac{2}{3}-\frac{1}{3}\times \frac{4}{5} &
	 b&= \frac{-2}{3} \times \left( \frac{1}{2}-\frac{1}{4} \right) \\
	 c&= \left( \frac{-2}{7}+\frac{5}{42} \right)\times \left( 5-\frac{3}{8} \right) & d&= \frac{\frac{2}{5}+\frac{-3}{4}}{2+(-2)\times \frac{-7}{4}} \\
	 e&=\left(1+\frac{1}{2}\right) \div \frac{2}{3} &
	f&=\left(\frac{2}{3}+\frac{1}{2}\right) \times \frac{6}{25} \\
	g&=\frac{-1+\frac{33}{11}}{\frac{2}{7}}
	\end{align*}
}{exe:Samoyeau5}{
	sol1
}

\exe{difficulty=0}{
	En décembre pour les fêtes, M. Marchand dit avoir vendu les quatre cinquièmes de sa marchandise. En janvier, pendant les soldes, il a encore vendu les trois quarts de ce qu'il restait.
	\begin{enumerate}
		\item
		Quelle fraction de sa marchandise a-t-il vendu en tout?
		\item  
		La valeur totale de sa marchandise est de 262 000 euros. Quelle somme représente sa vente globale?
	\end{enumerate}

}{exe:Samoyeau6}{
	sol1
}

\exe{difficulty=0}{
	Une usine française de voitures vend $\frac{1}{3}$ de ses voitures en Allemagne, $\frac{1}{4}$ de ce qui reste en Espagne et le reste en France. Quelle proportion de voitures de cette usine reste en France?

}{exe:Samoyeau7}{
	sol1
}

\exe{difficulty=0}{
	Amaya et Mariza ont deux tablettes de chocolat, une au chocolat noir et une au chocolat au lait qui sont de même proportion. 
	Amaya a mangé $\frac{1}{4}$ des $\frac{2}{3}$ de la tablette au chocolat au lait. Mariza a mangé la moitié des $\frac{3}{4}$ de la tablette de chocolat noir. Laquelle a mangé le plus de chocolat? 

}{exe:Samoyeau8}{
	sol1
}

\exe{difficulty=0}{
	Lors d'une course, Pierre fait $\frac{1}{3}$ de la distance totale parcourue en vélo, $\frac{5}{12}$ de la distance totale en courant, $\frac{1}{6}$ de la distance totale à la nage et le reste en kayak. 
	\begin{enumerate}
		\item Compléter la phrase suivante : Pierre a parcouru .... de la distance totale en kayak.
		\item Cette course fait 24 kilomètres. Quelle est la distance parcourue en kayak? 
	\end{enumerate} 

}{exe:Samoyeau9}{
	sol1
}

\exe{difficulty=0}{
	Calculer la fraction suivante et simplifier au maximum.
	\[ 
	\beta=\frac{ \frac{2}{3} \times \frac{6}{12} + \frac{-3}{7} - 4 }{ \frac{\frac{1}{2}-3}{8\times \frac{-1}{2}} + \frac{2}{5}-\frac{1}{4}}
	\]

}{exe:Samoyeau10}{
	sol1
}






\newpage
\addcontentsline{toc}{section}{Solutions}
\fancyhead[C]{\textbf{Chapitre \thechapter~— Solutions}}

\shipoutAnswer



